\documentclass[11pt,a4paper]{article}

\usepackage{iftex}
\RequirePDFTeX
\usepackage{cmap}
\usepackage[T2A]{fontenc}
\usepackage[utf8]{inputenc}
\usepackage[english,russian]{babel}
\usepackage{PTSerif}
\usepackage{amsmath,amssymb,microtype}
\usepackage[a4paper,left=23mm,right=23mm,top=20mm,bottom=21mm]{geometry}
\usepackage{xcolor,enumitem,fancyhdr,hyperref}

\definecolor{accent}{HTML}{187AA5}
\definecolor{ink}{HTML}{302B28}
\definecolor{muted}{HTML}{6D6560}
\definecolor{tint}{HTML}{EDF4F6}
\definecolor{rulecolor}{HTML}{A9CADA}
\hypersetup{unicode=true,colorlinks=true,urlcolor=accent,linkcolor=accent,
  pdftitle={Домашнее задание 1. Вековые возмущения от J2},
  pdfauthor={Чиняев Владимир Николаевич}}
\urlstyle{same}
\color{ink}
\setlength{\parindent}{0pt}
\setlength{\parskip}{5pt}
\linespread{1.08}
\setlist[enumerate]{leftmargin=*,label=\textcolor{accent}{\arabic*.},
  itemsep=5pt,topsep=5pt,parsep=0pt}
\setlist[itemize]{leftmargin=*,itemsep=3pt,topsep=3pt,parsep=0pt}
\setlength{\emergencystretch}{1em}
\pagestyle{fancy}
\fancyhf{}
\renewcommand{\headrulewidth}{0pt}
\renewcommand{\footrulewidth}{0.4pt}
\renewcommand{\footrule}{\hbox to\headwidth{\color{rulecolor}\leaders\hrule height\footrulewidth\hfill}}
\fancyfoot[L]{\footnotesize\color{muted}Введение в маневрирование КА\quad /\quad ДЗ 1}
\fancyfoot[R]{\footnotesize\color{muted}\thepage}

\newcommand{\taskheading}[2]{%
  \par\vspace{5pt}{\large\bfseries\color{accent}Задание #1. #2}\par\smallskip}

\begin{document}

{\small\color{accent}ВВЕДЕНИЕ В МАНЕВРИРОВАНИЕ КА}\par
\vspace{3pt}
{\huge Домашнее задание 1}\par
{\large Вековые возмущения от $J_2$}\par
{\small\color{muted}К лекции 2}
\par{\small\bfseries\color{accent}Срок сдачи: 21 сентября 2026 года, 23:59}
\vspace{4pt}
{\color{rulecolor}\hrule}
\vspace{5pt}

В заданиях рассматриваются вековые изменения средних орбитальных элементов,
обусловленные второй зональной гармоникой $J_2$ гравитационного поля Земли.
Расчёты выполняются в приближении первого порядка по $J_2$.
Остальные возмущения не учитываются.

Оформите решение обоих заданий в Jupyter Notebook \texttt{hw\_1.ipynb}.
Дополните предоставленный файл ячейками с расчётами, графиками,
аналитическими выводами и ответами на вопросы.

\begingroup
\setlength{\fboxsep}{9pt}
\colorbox{tint}{\begin{minipage}{\dimexpr\linewidth-2\fboxsep\relax}
\textbf{Константы и обозначения}
\[
\mu=398600{,}4418\ \text{км}^3/\text{с}^2,\qquad
R_\oplus=6378{,}1363\ \text{км},\qquad
J_2=1{,}08262668\cdot10^{-3}.
\]
\vspace{-6pt}
\[
n=\sqrt{\frac{\mu}{a^3}},\qquad p=a(1-e^2),\qquad
\dot\Omega=-\frac{3}{2}J_2n\left(\frac{R_\oplus}{p}\right)^2\cos i.
\]
При использовании этих единиц $n$ и $\dot\Omega$ выражены в рад/с.
Одни сутки равны $86400$~с. Аргументы тригонометрических функций задавайте в радианах.
\end{minipage}}
\endgroup

\taskheading{1}{Солнечносинхронная орбита}

В рассматриваемом приближении условие солнечносинхронности состоит в том,
что долгота восходящего узла увеличивается со средней угловой скоростью
годичного движения Солнца:
\[
\dot\Omega=\dot\Omega_{\text{ССО}}=+0{,}9856^\circ/\text{сут}.
\]
Такой выбор прецессии обеспечивает приблизительное постоянство местного
среднего солнечного времени прохождения восходящего узла.

Рассмотрите круговые орбиты ($e=0$): $a=R_\oplus+h$,
где $h$ --- высота над сферой радиуса $R_\oplus$.

\begin{enumerate}
\item Из условия солнечносинхронности получите выражение для наклонения $i(h)$.
Укажите условие существования действительного решения.
Определите предельную высоту круговой солнечносинхронной орбиты в принятой модели.
Можно ли обеспечить заданную скорость прецессии узла выше этой границы,
изменяя только наклонение? Объясните причину ограничения.
Какому наклонению соответствует сама граница и определён ли при нём восходящий узел?
\item Постройте график $i(h)$ на интервале $300\leq h\leq1500$~км.
По горизонтальной оси отложите $h$ в километрах, по вертикальной --- $i$ в градусах.
\item Вычислите наклонения для высот $400$, $700$ и $1200$~км.
Результаты представьте в таблице с точностью до $0{,}01^\circ$.
\item Объясните, почему полученные орбиты являются обратными
и как изменяется требуемое наклонение при увеличении высоты.
Обоснуйте ответ по формуле вековой скорости узла.
\end{enumerate}

\newpage

{\small\color{muted}Домашнее задание 1\hfill К лекции 2}
\taskheading{2}{Прецессия плоскости орбиты и линии апсид}

Заданы большая полуось $a=10000$~км и эксцентриситет $e=0{,}3$.
Рассматриваются орбиты с наклонениями $0<i<180^\circ$.
Вековая скорость аргумента перицентра в первом приближении по $J_2$ равна
\[
\dot\omega=\frac{3}{4}J_2n\left(\frac{R_\oplus}{p}\right)^2
\bigl(5\cos^2 i-1\bigr).
\]
Долгота восходящего узла $\Omega$ характеризует ориентацию плоскости орбиты,
а аргумент перицентра $\omega$ отсчитывается в плоскости орбиты от восходящего узла.

\begin{enumerate}
\item Постройте на одном графике зависимости $\dot\Omega(i)$ и $\dot\omega(i)$
для $5^\circ\leq i\leq175^\circ$.
Наклонение выразите в градусах, скорости --- в градусах за сутки.
Добавьте легенду и горизонтальную линию нулевой скорости.
\item Аналитически найдите все наклонения в интервале $0<i<180^\circ$,
при которых $\dot\Omega=0$ или $\dot\omega=0$
(наклонения, при которых $\dot\omega=0$, называются критическими).
Отметьте найденные значения на графике.
Определите, существует ли наклонение, при котором обе скорости равны нулю одновременно.
\item С помощью файла \texttt{hw\_1.ipynb} исследуйте пространственную эволюцию
трёх орбит: с наклонениями $50^\circ$, $90^\circ$ и первым из найденных
критических наклонений ($\dot\omega=0$, $i<90^\circ$).
Начальные углы равны $\Omega_0=20^\circ$, $\omega_0=40^\circ$.
Рассмотрите эпохи $t=0$, $90$ и $180$~сут.
\item Для полярной орбиты и орбиты с критическим наклонением объясните,
как изменяются плоскость орбиты, линия узлов и направление на перицентр.
Следует ли из равенства $\dot\omega=0$ постоянство направления на перицентр
в инерциальной системе координат? Обоснуйте ответ.
\end{enumerate}

\vspace{4pt}
\textbf{Работа с пространственной визуализацией}

Откройте \texttt{hw\_1.ipynb} в среде Jupyter и выполните ячейки по порядку.
В коде можно не разбираться совсем --- просто посмотрите на картинки и поймите из них что-то.
Рисунки для эпох $0$, $90$ и $180$~сут выводятся автоматически.
При наличии интерактивных элементов управления можно также выбирать промежуточные эпохи.
Используйте рисунки для ответов на вопросы задания~2.

Оси $X,Y,Z$ задают геоцентрическую инерциальную систему координат,
плоскость $XY$ является экваториальной. В модели $a,e,i$ постоянны,
а углы изменяются по законам
\[
\Omega(t)=\Omega_0+\dot\Omega t,\qquad
\omega(t)=\omega_0+\dot\omega t.
\]
Каждый эллипс изображает орбиту в средних элементах на выбранную эпоху.
Последовательность эллипсов не является траекторией движения КА во времени.

\vspace{5pt}
{\color{rulecolor}\hrule}
\vspace{5pt}
\textbf{Что представить}

Пришлите один файл \texttt{hw\_1.ipynb} с полным решением обоих заданий:
выводом выражения $i(h)$, таблицей наклонений, двумя графиками,
аналитическим определением нулей вековых скоростей и обоснованными ответами
на все вопросы. Формулы и пояснения оформите в текстовых ячейках,
расчёты и построение графиков --- в кодовых.
Перед отправкой выполните все ячейки по порядку и сохраните ноутбук
с результатами вычислений и графиками.

\end{document}
